\documentclass{article}
\usepackage[utf8]{inputenc}    
\usepackage[T1]{fontenc}       
\usepackage{lmodern}           
\usepackage{hyperref}
\hypersetup{
    colorlinks=true,       % false: boxed links; true: colored links
    linkcolor=blue,        % color of internal links (sections, etc.)
    urlcolor=cyan,         % color of external links
    linktoc=all            % 'all' makes both text and page number clickable
}
\usepackage{amsmath}   
\usepackage{amssymb}   
\usepackage{geometry}  
\usepackage{enumerate} 
\usepackage{xcolor}  
\usepackage{amsthm}
\usepackage{pdfpages}
\usepackage{mathrsfs}
\newtheorem{theorem}{Theorem}[section]
% \newtheorem{lemma}[theorem]{Lemma}
\newtheorem{corollary}[theorem]{Corollary}
% \newtheorem{definition}[theorem]{Definition}
\newtheorem{example}[theorem]{Example}
\usepackage{listings}  
\usepackage{tikz-cd}
\usepackage{forest}     
\usepackage{xcolor}
\usepackage{tcolorbox}
\tcbuselibrary{theorems}

% remove/comment out the old line:
% \newtheorem{definition}[theorem]{Definition}

\usepackage{xcolor}
\usepackage{tcolorbox}
\tcbuselibrary{breakable}
\newtheorem{definition}[theorem]{Definition}
\let\olddefinition\definition
\let\endolddefinition\enddefinition
\renewenvironment{definition}
  {\begin{tcolorbox}[
      colback=red!10,
      colframe=red!50!black,
      boxrule=0.6pt,
      arc=2pt,
      left=6pt,right=6pt,top=4pt,bottom=4pt,
      breakable
    ]\olddefinition}
  {\endolddefinition\end{tcolorbox}}

\newtheorem{lemma}[theorem]{Lemma}
\let\oldlemma\lemma
\let\endoldlemma\endlemma
\renewenvironment{lemma}
  {\begin{tcolorbox}[
      colback=green!8,
      colframe=green!50!black,
      boxrule=0.6pt,
      arc=2pt,
      left=6pt,right=6pt,top=4pt,bottom=4pt
    ]\oldlemma}
  {\endoldlemma\end{tcolorbox}}

  \let\oldtheorem\theorem
\let\endoldtheorem\endtheorem
\renewenvironment{theorem}
  {\begin{tcolorbox}[
      colback=yellow!6,
      colframe=yellow!60!black,
      boxrule=0.6pt,
      arc=2pt,
      left=6pt,right=6pt,top=4pt,bottom=4pt,
      parbox=false
    ]\oldtheorem}
  {\endoldtheorem\end{tcolorbox}}
\usepackage{bussproofs}
\usetikzlibrary{arrows.meta}
\usetikzlibrary{arrows.meta,decorations.pathreplacing,calc}
\lstset{frame=tb,
  language=C,
  aboveskip=3mm,
  belowskip=3mm,
  showstringspaces=false,   
  columns=flexible,
  basicstyle={\small\ttfamily},
  numbers=none,
  numberstyle=\tiny\color{gray},
  keywordstyle=\color{blue},
  commentstyle=\color{brown},
  stringstyle=\color{orange},
  breaklines=true,
  breakatwhitespace=true,
  tabsize=3
}
\geometry{top=0.5in, bottom=0.5in, left=1in, right=1in}
% ctrl shift D to toggle night 
\begin{document}

\title{HW3}
\author{Wang Xiyu A0282500R}
\date{}
\maketitle
\newpage
\section*{Hilbert-style Proof System}
The Hilbert-style (or Frege) proof system for propositional logic comprises of 3 axiom schema and the modus ponens rule of inference

\begin{itemize}
    \item \textbf{Axiom Schema 1:} $\varphi \rightarrow (\psi \rightarrow \varphi)$ 
    \item \textbf{Axiom Schema 2:} $(\varphi \rightarrow (\psi \rightarrow \rho)) \rightarrow ((\varphi \rightarrow \psi) \rightarrow (\varphi \rightarrow \rho))$ 
    \item \textbf{Axiom Schema 3:} $((\varphi \rightarrow \perp) \rightarrow \perp) \rightarrow \varphi$
    \item \textbf{Modus Ponens:} From $\varphi$ and $\varphi \rightarrow \psi$, infer $\psi$
\end{itemize}

\section*{Question 1}
Let $\Gamma, \Gamma'$ be sets of propositional logic formulae and let $\varphi$ be a formula. Show that if $\Gamma \vdash \varphi$ and if $\Gamma \subseteq \Gamma'$, then $\Gamma' \vdash \varphi$
\begin{proof}
    $\Gamma \vdash \varphi \implies \exists $ a proof for $\phi, \langle \phi_1, \phi_2, \dots, \phi_n\rangle, \phi_n = \varphi$
    \subsubsection*{Base case: $n = 1$}
    \begin{itemize}
        \item Case 1: $\varphi \in \Gamma \implies \varphi \in \Gamma' \implies \Gamma' \vdash \varphi$
        \item Case 2: $\varphi$ is from Axiom 1 or 2 or 3, $\vdash \varphi \implies \Gamma' \vdash \varphi$
    \end{itemize}
    \subsubsection*{Inductive case: Suppose $\forall i < n, \Gamma \vdash \phi_i \implies \Gamma' \vdash \phi_i$}
    \begin{itemize}
        \item Case 1: $\varphi$ is from axiom 1, 2, 3, trivial
        \item Case 2: $\varphi$ is from Modus Ponens, $\exists i, j < n, \phi_i \equiv \phi_j \rightarrow \varphi, \Gamma \vdash \phi_i, \Gamma \vdash \phi_j$ 
        \begin{prooftree}
            \AxiomC{$\phi_j \to \varphi$}
            \AxiomC{$\phi_j$}
            \BinaryInfC{$\varphi$}
        \end{prooftree}
        From indeuctive hypothesis, $\Gamma' \vdash \phi_i, \Gamma' \vdash \phi_j, \Gamma' \vdash \varphi$
    \end{itemize}
\end{proof}
\section*{Question 2}

\begin{theorem}[Deduction Theorem]
    Let $\Gamma$ be a set of propositional logic formulae and $\varphi, \psi$ be propositional logic formulae. $\Gamma \vdash (\varphi \rightarrow \psi)$ iff $\Gamma \cup \{\varphi\} \vdash \psi$.
\end{theorem}

\begin{proof}
    \subsubsection*{($\implies$)} 
    \begin{enumerate}
        \item $\Gamma \vdash \varphi \rightarrow \psi$ (Given)
        \item $\Gamma \cup \{\varphi\} \vdash \varphi \rightarrow \psi$ (Monotonicity)
        \item $\Gamma \cup \{\varphi\} \vdash \varphi$ (Assumption)
        \item $\Gamma \cup \{\varphi\} \vdash \psi$ (MP 2, 3)
    \end{enumerate}

    \subsubsection*{($\impliedby$)} 
    Let $\pi = \psi_1, \dots, \psi_n$ be a proof of $\psi$ from $\Gamma \cup \{\varphi\}$. Assume $\forall j < i, \Gamma \vdash \varphi \rightarrow \psi_j$.

    \paragraph{Case 1: $\psi_i \in \Gamma$ or $\psi_i$ is an Axiom}
    \begin{enumerate}
        \item $\Gamma \vdash \psi_i$ (Assumption/Axiom)
        \item $\Gamma \vdash \psi_i \rightarrow (\varphi \rightarrow \psi_i)$ (Ax 1)
        \item $\Gamma \vdash \varphi \rightarrow \psi_i$ (MP 1, 2)
    \end{enumerate}

    \paragraph{Case 2: $\psi_i = \varphi$}
    \begin{enumerate}
        \item $\Gamma \vdash (\varphi\rightarrow ((\varphi \rightarrow \varphi)\rightarrow \varphi)) \rightarrow ((\varphi\rightarrow (\varphi \rightarrow \varphi))\rightarrow (\varphi \rightarrow \varphi))$ (Ax 2)
        \item $\Gamma \vdash \varphi\rightarrow((\varphi\rightarrow \varphi)\rightarrow \varphi)$ (Ax 1)
        \item $\Gamma \vdash (\varphi\rightarrow(\varphi\rightarrow \varphi)) \rightarrow(\varphi \rightarrow \varphi)$ (MP 1, 2)
        \item $\Gamma \vdash \varphi \rightarrow(\varphi \rightarrow \varphi)$ (Ax 1)
        \item $\Gamma \vdash (\varphi \rightarrow \varphi)$ (MP 3, 4)
    \end{enumerate}

    \paragraph{Case 3: $\psi_i$ via MP from $\psi_j, \psi_k$ where $\psi_j = \psi_k \rightarrow \psi_i$}
    \begin{enumerate}
        \item $\Gamma \vdash \varphi \rightarrow (\psi_k \rightarrow \psi_i)$ (IH)
        \item $\Gamma \vdash (\varphi \rightarrow (\psi_k \rightarrow \psi_i)) \rightarrow ((\varphi \rightarrow \psi_k) \rightarrow (\varphi \rightarrow \psi_i))$ (Ax 2)
        \item $\Gamma \vdash (\varphi \rightarrow \psi_k) \rightarrow (\varphi \rightarrow \psi_i)$ (MP 1, 2)
        \item $\Gamma \vdash \varphi \rightarrow \psi_k$ (IH)
        \item $\Gamma \vdash \varphi \rightarrow \psi_i$ (MP 3, 4)
    \end{enumerate}

\end{proof}
\section*{Question 3}
\subsection*{(a) $\vdash \varphi \rightarrow \varphi$}
Done in Q2

\subsection*{(b) $\vdash (\psi \rightarrow \varphi) \rightarrow \varphi$, where $\psi = \varphi \rightarrow \perp$}
\begin{proof}
    $\equiv \{\psi \rightarrow \varphi\} \vdash \varphi$.
    \begin{enumerate}
        \item $(\varphi \rightarrow \perp) \rightarrow \varphi$ (Assumption)
        \item $((\varphi \rightarrow \perp) \rightarrow \perp) \rightarrow \varphi$ (Axiom 3)
        \item $((\varphi \rightarrow \perp) \rightarrow \varphi) \rightarrow (((\varphi \rightarrow \perp) \rightarrow (\varphi \rightarrow \perp)) \rightarrow ((\varphi \rightarrow \perp) \rightarrow \perp))$ (Axiom 2)
        \item $((\varphi \rightarrow \perp) \rightarrow (\varphi \rightarrow \perp)) \rightarrow ((\varphi \rightarrow \perp) \rightarrow \perp)$ (MP 1, 3)
        \item $(\varphi \rightarrow \perp) \rightarrow (\varphi \rightarrow \perp)$ (Part a)
        \item $(\varphi \rightarrow \perp) \rightarrow \perp$ (MP 4, 5)
        \item $\varphi$ (MP 2, 6)
    \end{enumerate}
    (Wikipedia: Peirce's law).
\end{proof}

\subsection*{(c) $\{\perp\} \vdash \varphi$}
\begin{proof}
    \begin{enumerate}
        \item $\perp$ (Assumption)
        \item $\perp \rightarrow ((\varphi \rightarrow \perp) \rightarrow \perp)$ (Axiom 1)
        \item $(\varphi \rightarrow \perp) \rightarrow \perp$ (MP 1, 2)
        \item $((\varphi \rightarrow \perp) \rightarrow \perp) \rightarrow \varphi$ (Axiom 3)
        \item $\varphi$ (MP 3, 4)
    \end{enumerate}
\end{proof}

\subsection*{(d) $\vdash \psi \rightarrow (\varphi \rightarrow \rho)$ where $\psi = \varphi \rightarrow \perp$}
\begin{proof}
    $\equiv \{\varphi \rightarrow \perp, \varphi\} \vdash \rho$.
    \begin{enumerate}
        \item $\varphi \rightarrow \perp$ (Assumption)
        \item $\varphi$ (Assumption)
        \item $\perp$ (MP 1, 2)
        \item $\perp \rightarrow ((\rho \rightarrow \perp) \rightarrow \perp)$ (Axiom 1)
        \item $(\rho \rightarrow \perp) \rightarrow \perp$ (MP 3, 4)
        \item $((\rho \rightarrow \perp) \rightarrow \perp) \rightarrow \rho$ (Axiom 3)
        \item $\rho$ (MP 5, 6)
    \end{enumerate}
\end{proof}

\section*{Theorem 2 (Soundness of Hilbert-style Proof System)}
Let $\Gamma$ be a set of propositional logic formulae and $\varphi$ be a propositional logic formula. If $\Gamma \vdash \varphi$, then $\Gamma \models \varphi$

\section*{Question 4}
Prove Theorem 2. \textit{Hint: Given a proof, inductively establish that every formula in the proof is a logical consequence of the assumptions.}
\begin{proof}
    Let $\pi=\phi_1,\phi_2,\dots,\phi_n$ be a proof of $\varphi$ from $\Gamma$, so $\phi_n=\varphi$
    \[
    \Gamma \models \phi_i
    \]
    $\forall 1\le i\le n$.
    
    \begin{enumerate}
        \item If $\phi_i \in \Gamma \implies \Gamma \models \phi_i$.
    
        \item If $\phi_i$ from an axiom, $\models \phi_i$
    
        \item Otherwise, $\phi_i$ is by MP from $\phi_j$ and $\phi_k=\phi_j\rightarrow \phi_i$, where $j,k<i$. By IH
        \[
        \Gamma \models \phi_j, \qquad
        \Gamma \models \phi_j\rightarrow \phi_i.
        \]
        Let $v\models \Gamma, v\models \phi_j, v\models \phi_j\rightarrow \phi_i, v\models (\neg \phi_j \lor \phi_i), v\models \phi_i$
    \end{enumerate}
    
    Therefore $\Gamma \models \phi_n = \varphi$. Thus $\Gamma \vdash \varphi$ implies $\Gamma \models \varphi$.
    \end{proof}
\section*{Theorem 3 (Completeness of Hilbert-style Proof System)}
Let $\Gamma$ be a set of propositional logic formulae and $\varphi$ be a propositional logic formula. If $\Gamma \models \varphi$, then $\Gamma \vdash \varphi$

\section*{Theorem 4 (Weak Completeness of Hilbert-style Proof System)}
Let $\varphi$ be a propositional logic formula. If $\models \varphi$, then $\vdash \varphi$.
\begin{definition}[Inconsistent set]
    $\Gamma$ is inconsistent $\iff$
    \[\exists \varphi, \Gamma \vdash \varphi, \Gamma \vdash (\varphi \rightarrow \bot)\]
\end{definition}
\section*{Question 5}
Prove that two notions of a set being inconsistent are identical. That is, prove that the following two statements are equivalent
\begin{enumerate}
    \item There is a formula $\varphi$ such that $\Gamma \vdash \varphi$ and $\Gamma \vdash (\varphi \rightarrow \perp)$
    \item For every formula $\rho$, $\Gamma \vdash \rho$
\end{enumerate}
\begin{proof}
    $(1 \Rightarrow 2)$: If $\Gamma \vdash \varphi$ and $\Gamma \vdash \varphi \rightarrow \bot$, then by MP $\Gamma \vdash \bot$; since $\{\bot\} \vdash \rho$ for every formula $\rho$, by monotonicity we get $\Gamma \vdash \rho$ for every formula $\rho$.
    
    $(2 \Rightarrow 1)$: for any formula $\varphi$, taking $\rho=\varphi$ and $\rho=\varphi \rightarrow \bot$ gives $\Gamma \vdash \varphi$ and $\Gamma \vdash \varphi \rightarrow \bot$.
    \end{proof}
\section*{Lemma 1}
Let $\Gamma$ be a set of formulae and $\varphi$ be a formula. If $\Gamma \cup \{\varphi \rightarrow \perp\}$ is inconsistent, then $\Gamma \vdash \varphi$

\section*{Lemma 2}
Let $\Gamma$ be a set of formulae and $\varphi$ be a formula. If $\Gamma \cup \{\varphi\}$ is inconsistent, then $\Gamma \vdash (\varphi \rightarrow \perp)$

\section*{Lemma 3}
Let $\Gamma$ be a set of formulae. If $\Gamma$ is inconsistent, then there is a finite subset $\Gamma' \subseteq_{fin} \Gamma$ of $\Gamma$ such that $\Gamma'$ is also inconsistent

\section*{Question 6}
\begin{itemize}
    \item (a) Prove Lemma 1
        \begin{proof}
            From deduction thm, $\Gamma \vdash (\varphi \rightarrow \perp) \rightarrow \psi$ and $\Gamma \vdash (\varphi \rightarrow \perp) \rightarrow (\psi \rightarrow \perp)$.
            \begin{enumerate}
                \item $\Gamma \vdash (\varphi \rightarrow \perp) \rightarrow (\psi \rightarrow \perp)$ 
                \item $\Gamma \vdash ((\varphi \rightarrow \perp) \rightarrow (\psi \rightarrow \perp)) \rightarrow (((\varphi \rightarrow \perp) \rightarrow \psi) \rightarrow ((\varphi \rightarrow \perp) \rightarrow \perp))$ (Ax 2)
                \item $\Gamma \vdash ((\varphi \rightarrow \perp) \rightarrow \psi) \rightarrow ((\varphi \rightarrow \perp) \rightarrow \perp)$ (MP 1, 2)
                \item $\Gamma \vdash (\varphi \rightarrow \perp) \rightarrow \psi$
                \item $\Gamma \vdash (\varphi \rightarrow \perp) \rightarrow \perp$ (MP 3, 4)
                \item $\Gamma \vdash ((\varphi \rightarrow \perp) \rightarrow \perp) \rightarrow \varphi$ (Ax 3)
                \item $\Gamma \vdash \varphi$ (MP 5, 6)
            \end{enumerate}
        \end{proof}
    \item (b) Prove Lemma 2
    \begin{proof}
        By definition, $\Gamma \cup \{\varphi\} \vdash \psi$ and $\Gamma \cup \{\varphi\} \vdash (\psi \rightarrow \perp)$ for some $\psi$. 
        By the Deduction:
        \begin{enumerate}
            \item $\Gamma \vdash \varphi \rightarrow \psi$
            \item $\Gamma \vdash \varphi \rightarrow (\psi \rightarrow \perp)$
            \item $\Gamma \vdash (\varphi \rightarrow (\psi \rightarrow \perp)) \rightarrow ((\varphi \rightarrow \psi) \rightarrow (\varphi \rightarrow \perp))$ (Ax 2)
            \item $\Gamma \vdash (\varphi \rightarrow \psi) \rightarrow (\varphi \rightarrow \perp)$ (MP 2, 3)
            \item $\Gamma \vdash \varphi \rightarrow \perp$ (MP 1, 4)
        \end{enumerate}
    \end{proof}
    \item (c) Let $\Gamma$ be a set of formulae and $\varphi$ be a formula. Show that, if $\Gamma$ is consistent, then at least one of $\Gamma \cup \{\varphi\}$ or $\Gamma \cup \{\varphi \rightarrow \perp\}$ is consistent
    \begin{proof}
        Suppos otherwise, $\Gamma$ is consistent, $\Gamma \cup \{\varphi\}$ and $\Gamma \cup \{\varphi \rightarrow \perp\}$ are inconsistent.
        Since $\Gamma \cup \{\varphi\}$ is inconsistent, by Lemma 2
        \[
        \Gamma \vdash \varphi \rightarrow \perp.
        \]
        Since $\Gamma \cup \{\varphi \rightarrow \perp\}$ is inconsistent, by Lemma 1 we have
        \[
        \Gamma \vdash \varphi.
        \]
        \[
        \Gamma \vdash \perp.
        \]
        By explosion, pick an arbitrary formula $\psi, 
        \Gamma \vdash \bot \vdash \psi, \Gamma \vdash \bot \vdash \psi \rightarrow \bot, \Gamma$ is inconsistent, contradiction.
        \end{proof}
    \item (d) Let $\Gamma, \Gamma'$ be sets of formulae such that $\Gamma \subseteq \Gamma'$. If $\Gamma'$ is consistent, then so is $\Gamma$
    \begin{proof}
        Suppose othrewise, $\Gamma$ is inconsistent,
        \[
        \Gamma \vdash \perp.
        \]
        $\Gamma \subseteq \Gamma'$, by monotonicity of derivability we have
        \[
        \Gamma' \vdash \perp.
        \]
        Contradiction.
        \end{proof}
    \item (e)
    \begin{proof}
        $\Gamma $ is inconsistent, $\exists \varphi,  \Gamma \vdash \varphi, \Gamma \vdash (\varphi \rightarrow \bot)$
        there exists a finite proof for both $\varphi$ and $\varphi \rightarrow \bot$. Consider 
        \begin{itemize}
            \item $\pi_1 = \langle\phi_1, \phi_2, \dots, \phi_n\rangle, \phi_n = \varphi$
            \item $\pi_2 = \langle\psi_1, \psi_2, \dots, \psi_m\rangle, \psi_m = \varphi \rightarrow \bot$
        \end{itemize}
        We construct $\Gamma'$ as such: $\forall i < n, j < m, $ if $\phi_i \in \Gamma, \phi_i \in \Gamma'$; if $\psi_j \in \Gamma, \psi_j \in \Gamma'$ 
        Therefore $\forall \eta \in \Gamma', \eta \in \Gamma, \Gamma' \subseteq \Gamma$; $|\Gamma'| \leq n + m, \Gamma'$ is a finite subset of $\Gamma$.
        Since $pi_1$ is a proof for $\phi$, $\forall \phi \in \pi_1, \phi \in \Gamma \implies \phi \in \Gamma'$; We can easily show that same proof can be 
        constructed from $\Gamma'$ by induction. Same case for $\Gamma'\vdash \varphi \rightarrow \bot$. Therefore $\Gamma'$ is inconsistent.       
    \end{proof}
\end{itemize}
\begin{definition}
    \begin{itemize}
        \item $\Delta_0 = \{\varphi \rightarrow \bot\}$
        \item $\Delta_{i+1} = \begin{cases}
            \Delta_i \cup \{\varphi_i\} \text{ if } \Delta_i \cup \{\varphi_i\} \text{ is consistent}\\
            \Delta_i \cup \{\varphi_i \rightarrow \bot\} \text{ otherwise}
        \end{cases}$
    \end{itemize}
    \[\Delta = \bigcup_{i\geq 0} \Delta_i\]
\end{definition}
\begin{lemma}
    $\Gamma$ is consistent $\implies \Gamma \cup \{\varphi\}$ or $\Gamma \cup \{\varphi \rightarrow \bot\}$ is consistent.
\end{lemma}

\section*{Question 7}
\begin{itemize}
    \item (a) Prove, using induction, that $\Delta_i$ is consistent, for every $i$
    \begin{proof}
        \subsubsection*{Base case: $\Delta_0$}
        Suppose otherwise, $\Delta_0=\{\varphi \rightarrow \bot\}$ is inconsistent. Then by Lemma 1, 
        \[
        \vdash \varphi.
        \]
        Contradiction.
        \subsubsection*{Inductive case: Suppose $\forall j \leq i, \Delta_j$ is consistent}
        By lemma above, $\Delta_i$ is consistent $\implies \Delta_i \cup \{\varphi_i\}$ is consistent or $\Delta_i \cup \{\varphi_i \rightarrow \bot\}$
        is consistent $\implies \Delta_{i+1}$ is consistent
    \end{proof}

    \item (b) Prove that $\Delta$ is consistent
    \begin{proof}
        Suppose otherwise, $\Delta$ is inconsistent, by Lemma 3, $\exists \Delta' \subseteq_{fin} \Delta, \Delta'$ is inconsistent. Since $\Delta'$ is finite, let the highest ranked formula in $\Delta'$
        be $n, \forall \varphi \in \Delta', \varphi \in \Delta_n$. From b), $\Delta_n$ is consistent, by 6d), $\Delta'$ is consistent, contradiction.
    \end{proof}
    \item (c) Prove that $\Delta$ is negation-complete, i.e., for every formula $\psi$, we have $(\psi \rightarrow \perp) \in \Delta$ iff $(\psi \notin \Delta)$
    \begin{proof}
        \subsubsection*{$\implies$}
        Let $\psi = \varphi_{n+1}, \Delta_{n}$ is consistent. $\Delta_{n+1} = \Delta_n \cup \{\psi\}$ if $\Delta_n \cup \{\psi\}$ is consistnet. Suppose $\psi \rightarrow \bot\in \Delta$, 
        $\Delta \vdash \psi, \Delta \vdash \psi \rightarrow \bot \implies \Delta$ is inconsitent, contradiction. Therefore $\psi \in \Delta \implies \psi \rightarrow \bot \notin \Delta$ 
        \subsubsection*{$\impliedby$}
        Suppose otherwise, $\psi \rightarrow \bot \in \Delta, \psi \in \Delta, \Delta \vdash \psi, \Delta \vdash \psi \rightarrow \bot \implies \Delta$ is inconsistent, contradiction.
    \end{proof}
    \item (d) Use the fact that $\Delta$ is negation-complete to show that for every formula $\psi$, $\psi \in \Delta$ iff $\Delta \vdash \psi$
    \begin{proof}
        \subsubsection*{$\implies$}
        trivial by definition
        \subsubsection*{$\impliedby$}
        Consider $\Delta \vdash \psi$. Suppose $\psi \notin \Delta, \psi \rightarrow \bot \in \Delta, 
        \Delta \vdash \psi \rightarrow \bot$. 
        Since $\Delta \vdash \psi$, by MP, $\Delta \vdash \bot$, 
        $\Delta$ is inconsistent, contradiction. Therefore $\psi \in \Delta$.
        \end{proof}
\end{itemize}













\section*{Question 8}
\begin{itemize}
    \item (a) Show that for any formula $\psi$, $v \models \psi$ iff $\psi \in \Delta$
    \begin{proof}
        \subsubsection*{Base case:}
        \begin{enumerate}
            \item $\psi \in P$: by definition, $v(p) = true \iff p \in \Delta$
            \item $\psi = \bot$: always $v(\bot) = false, \bot \notin \Delta \implies v\models \bot \iff \bot \in \Delta$ vacuously
        \end{enumerate}
        \subsubsection*{Inductive case: $\psi = a \rightarrow b$}
        Suppose 
        \begin{itemize}
            \item $v\models a \iff a \in \Delta$
            \item $v\models b \iff b \in \Delta$
        \end{itemize}
        \subsubsection*{Case $\implies$}
        Suppose $v\models a\rightarrow b$, $a \rightarrow b \notin \Delta, (a \rightarrow b) \rightarrow \bot \in \Delta$, 
        $\Delta \vdash (a\rightarrow b)\rightarrow \bot.$\\
        Now we show that $b \notin \Delta$
        \begin{itemize}
            \item Suppose otherwise, $b\in \Delta, \Delta \vdash b$, by ax 1, $\Delta \vdash b \rightarrow (a \rightarrow b) \implies \Delta \vdash a \rightarrow b$ by MP, contradiction.
        \end{itemize}
        Now we show that $a \in \Delta$
        \begin{itemize}
            \item Suppose otherwise, $a \not \in \Delta \implies a \rightarrow \bot \in \Delta$, from 
             \[\vdash (\varphi \rightarrow \bot)\rightarrow (\varphi \rightarrow \rho)\]
            $\Delta \vdash (a \rightarrow \bot) \rightarrow (a \rightarrow b)$. By MP, $\Delta \vdash a \rightarrow b$, contradiction.
        \end{itemize}
        Now with $b \notin \Delta$ and $a \in \Delta$, by IH, $v\models a$ and $v\not \models b \implies v\not \models a\rightarrow b$, contradiction.
        \subsubsection*{Case $\impliedby$}
        Suppose otherwise, $a\rightarrow b \in \Delta$, $v\models a$ and $v\not \models b$. 
        By IH, $a \in \Delta$ and $b \notin \Delta$, $b \rightarrow \bot \in \Delta\implies \Delta \vdash a, \Delta \vdash a \rightarrow b,$ by MP,
        $\Delta \vdash b \implies b \in \Delta$, contradiction.
    \end{proof}
    \item (b) Show that $v \models (\varphi \rightarrow \perp)$ and thus $\not\models \varphi$
    \begin{proof}
        $\{\varphi \rightarrow \bot\} \in \Delta \implies v \models (\varphi \rightarrow \bot) \implies v \not \models \varphi \implies \not \models \varphi$.
    \end{proof}
\end{itemize}

\section*{Question 9}
Prove Theorem 3 using Theorem 4, the Compactness theorem and the Deduction theorem
\begin{proof}
    $\Gamma \models \varphi \implies \Gamma \cup \{\varphi \rightarrow \bot\}$ is unsat, so by Compactness
    $
    \exists \Gamma_0 = \{\psi_1,\dots,\psi_n\} \subseteq_{\mathrm{fin}} \Gamma,
    $ such that
    $
    \Gamma_0 \cup \{\varphi \rightarrow \bot\}$
    is unsat, 
    \[ \{\psi_1, \dots, \psi_n, (\varphi \rightarrow \perp)\} \text{ is unsat} \implies 
    (\psi_1 \land \dots \land \psi_n \land (\varphi \rightarrow \perp)) \models \perp
    \]

    \[
    \models \psi_1 \rightarrow (\psi_2 \rightarrow \dots (\psi_n \rightarrow ((\varphi \rightarrow \bot)\rightarrow \bot))\dots)
    \]
    (https://en.wikipedia.org/wiki/Exportation\_(logic)) by Theorem 4
    \[
    \emptyset \vdash \psi_1 \rightarrow (\psi_2 \rightarrow \dots (\psi_n \rightarrow ((\varphi \rightarrow \bot)\rightarrow \bot))\dots)
    \]
    Apply deduction thm 
    \[\{\psi_1\}\vdash (\psi_2 \rightarrow \dots (\psi_n \rightarrow ((\varphi \rightarrow \bot)\rightarrow \bot))\dots)\]
    until we get
    \[
    \Gamma_0 \vdash (\varphi \rightarrow \bot)\rightarrow \bot.
    \]
    Using ax 3:
    \[
    \vdash ((\varphi \rightarrow \bot)\rightarrow \bot)\rightarrow \varphi,
    \]
    By MP
    \[
    \Gamma_0 \vdash \varphi.
    \]
    Since $\Gamma_0 \subseteq \Gamma$, by monotonicity,
    \[
    \Gamma \vdash \varphi.
    \]
    \end{proof}
\section*{Theorem 5 (Soundness of Resolution Proof System)}
Let $\Gamma$ be a set of clauses (over propositions $P$). If there is a resolution proof of $\Gamma$, then $\Gamma$ is unsatisfiable.

\section*{Question 10}
Prove Theorem 5, \textit{Hint: Consider a resolution proof $\pi = C_1, C_2, \dots, C_n$ of $\Gamma$. Use induction to establish that for every $1 \le i \le n$ if a valuation $v$ is such that $v \models \Gamma,$ then $v \models C_i$.}
\begin{proof}
    Let $\pi = C_1, C_2, \dots, C_n$ be a resolution proof of $\Gamma$, where $C_n = \emptyset$.
    
    We prove by induction on $i$ that for every $1 \leq i \leq n$, $v \models \Gamma, v \models C_i$.
    
    $i \in \{1,\dots,n\}$ and suppose $v \models \Gamma$.
    \begin{enumerate}
        \item $C_i \in \Gamma$, $v \models C_i$.
    
        \item $C_i$ is a resolvent of two earlier clauses $C_j$ and $C_k$ wrt some proposition $p$, where
        \[
        C_j = C \uplus \{p\}, \qquad C_k = D \uplus \{\neg p\},
        \]
        and
        $
        C_i = C \cup D.
        $
        By the IH, $j,k < i, v \models C_j \quad \text{and} \quad v \models C_k$
        We show that $v \models C_i$.
        Suppose, $v \not\models C_i = C \cup D$. Then $v$ satisfies neither any literal in $C$ nor any literal in $D$.
        Since $v \models C_j = C \uplus \{p\}$ and no literal in $C$ is satisfied by $v$, it must be that $v \models p$, i.e. $v(p)=\mathrm{true}$
        Since $v \models C_k = D \uplus \{\neg p\}$ and no literal in $D$ is satisfied by $v$, it must be that $v \models \neg p$, i.e. $v(p)=\mathrm{false}$,
    Contradicton
    \end{enumerate}
    \end{proof}
\section*{Theorem 6 (Completeness of Resolution Proof System)}
Let $\Gamma$ be a set of clauses (over propositions $P$). If $\Gamma$ is unsatisfiable, then there is a resolution proof of $\Gamma$

\section*{Question 11}
Let $\Gamma$ be an unsatisfiable set of clauses Show that, if $\Gamma$ is unsatisfiable, then there is a finite set $\Gamma_0 \subseteq_{fin} \Gamma$ such that $\Gamma_0$ is unsatisfiable.
\begin{proof}
    Suppose $\Gamma$ is unsat,by Compactness, $\exists \Gamma' \subseteq_{fin}\Gamma, \Gamma' $ is unsat (?)
    \end{proof}
\section*{Question 12}
Let $\Gamma_1, \Gamma_2$ be sets of clauses such that $\Gamma_1 \subseteq \Gamma_2$. Show that, if there is a resolution proof of $\Gamma_1$, then there is also a resolution proof of $\Gamma_2$
\begin{proof}
    Let $\pi = C_1, C_2, \dots, C_n$ be a resolution proof of $\Gamma_1$
    
    $\forall 1 \leq i \leq n$, $\pi$ is a resolution proof of $\Gamma_1$, either:
    \begin{enumerate}
        \item $C_i \in \Gamma_1$, or
        \item $\exists j, k < i$, $C_i$ is a resolvent of $C_j$ and $C_k$
    \end{enumerate}
    If $C_i \in \Gamma_1, C_i \in \Gamma_2$
    Otherwise, this fact is independent of the choice of the initial clause set, so the same derivation is valid in a proof from $\Gamma_2$.
    \end{proof}
\section*{Question 13}
    Next, we constructed a sequence of sets as follows. Here, we assume that the set of
    propositions occurring in $\Gamma_0$ can be enumerated as $p_1, p_2, \dots, p_k$.
    
    \[
    \Delta_0 = \Gamma_0
    \]
    
    \[
    \Delta_{i+1} =
    \Delta_i^{\emptyset}
    \;\cup\;
    \{\, C \cup D \mid
    C \uplus \{p_{i+1}\} \in \Delta_i^{+},
    \;
    D \uplus \{\neg p_{i+1}\} \in \Delta_i^{-}
    \,\}
    \]
    
    where
    
    \begin{itemize}
    \item $\Delta_i^{\emptyset}$ is the set of clauses in $\Delta_i$ that contain neither the literal $p_{i+1}$ nor $\neg p_{i+1}$,
    \item $\Delta_i^{+}$ is the set of clauses in $\Delta_i$ that contain the literal $p_{i+1}$ but not $\neg p_{i+1}$,
    \item $\Delta_i^{-}$ is the set of clauses in $\Delta_i$ that contain the literal $\neg p_{i+1}$ but not $p_{i+1}$.
    \end{itemize}
    
    Inductively show that $\Delta_i$ is unsat for every $0 \le i \le k$.
    Also show that $\Delta_i \ne \emptyset$ for every $0 \le i \le k$.


\begin{proof}
    \subsubsection*{Base case:} $i=0$.
    By construction, $\Delta_0 = \Gamma_0$. From 11, $\Gamma_0$ is a finite unsat subset of $\Gamma$
    \subsubsection*{Inductive step:}
    Suppose $\Delta_i$ is unsat and $\Delta_i \neq \emptyset$. We show $\Delta_{i+1}$ is unsat and nonempty.
    Suppose that $\exists v, v \models \Delta_{i+1}$, consider $v(p_{i+1})$
    \begin{itemize}
    \item If $v(p_{i+1}) = \text{true}$, then every clause in $\Delta_i^{+}$ is sat by $v$ due to $p_{i+1}$, and $\Delta_i^{\emptyset} \subseteq \Delta_{i+1}$ thus sat by $v$. 
    
    For every resolvent $C \cup D \in \Delta_{i+1}$ with
    \[
    C \uplus \{p_{i+1}\} \in \Delta_i^{+}, \quad
    D \uplus \{\neg p_{i+1}\} \in \Delta_i^{-},
    \]
    $v \models C \cup D$. Since $v(p_{i+1}) = \text{true}$, the clause $D \uplus \{\neg p_{i+1}\}$ must be satisfied by some literal in $D$, hence $v \models D$. Thus $v$ satisfies every clause in $\Delta_i^{-}$. Contradiction    
    \item case $v(p_{i+1}) = \text{false}$: similar
    \end{itemize}
\end{proof}
\section*{Question 14}
Inductively show that $\Delta_i$ does not contain the literals $\{p_1, \neg p_1, \dots, p_i, \neg p_i\}$ for every $1 \le i \le k$
\begin{proof}
    We prove by induction on $i$ that $\Delta_i$ does not contain any of the literals
    \[
    \{p_1, \neg p_1, \dots, p_i, \neg p_i\}
    \]
    for every $1 \leq i \leq k$.
    
    \subsubsection*{Base case:} $i=1$.
    
    By definition,
    \[
    \Delta_1
    =
    \Delta_0^{\emptyset}
    \cup
    \{\, C \cup D \mid C \uplus \{p_1\} \in \Delta_0^{+},\; D \uplus \{\neg p_1\} \in \Delta_0^{-}\,\}.
    \]
    Every clause in $\Delta_0^{\emptyset}$ contains neither $p_1$ nor $\neg p_1$, and every resolvent $C \cup D$ is obtained by removing $p_1$ and $\neg p_1$. Hence no clause in $\Delta_1$ contains $p_1$ or $\neg p_1$.
    \subsubsection*{Inductive step:}

    \end{proof}
\section*{Question 15}
Finish the proof of Theorem 6 by explicitly constructing a resolution proof for $\Gamma$ and establishing its correctness
\begin{proof}
Let $\Gamma$ be unsat. By 11, $\exists 
\Gamma_0 \subseteq_{{fin}} \Gamma
$
such that $\Gamma_0$ is unsat. Let $p_1,\dots,p_k$ be all propositions occured in $\Gamma_0$, and
$
\Delta_0,\Delta_1,\dots,\Delta_k
$
as in the construction. $\Delta_k$ is unsat, and by 14, $\Delta_k$ does not contain 
$
p_1,\neg p_1,\dots,p_k,\neg p_k.
$
therefore every clause in $\Delta_k$ contains no literal,thus every clause in $\Delta_k$ is $\emptyset$. Since $\Delta_k \neq \emptyset$,
$
\emptyset \in \Delta_k
$

\subsubsection*{Base case:}
$i=0$, $\Delta_0=\Gamma_0$
\subsubsection*{Inductive step:}
Suppose every clause in $\Delta_i$ has a finite resolution proof from $\Gamma_0$. Let $E \in \Delta_{i+1}$. By definition ,

\begin{enumerate}
    \item $E \in \Delta_i^{\emptyset}$, hence $E \in \Delta_i$, so by IH $E$ has a finite resolution proof from $\Gamma_0$; or
    \item $E = C \cup D$, where
    \[
    C \uplus \{p_{i+1}\} \in \Delta_i^{+}
    \qquad\text{and}\qquad
    D \uplus \{\neg p_{i+1}\} \in \Delta_i^{-}.
    \]
    Then $E$ is the resolvent of these 2
    \[
    C \uplus \{p_{i+1}\}
    \quad\text{and}\quad
    D \uplus \{\neg p_{i+1}\}.
    \]
    By IH, both clauses have a finite resolution proof from $\Gamma_0$, thus the resolution proof for $E$ is also finite
\end{enumerate}

$\Gamma_0 \subseteq \Gamma$, by 12 the same proof is also a resolution proof fo $\Gamma$.

Therefore $\Gamma$ is unsat, then there is a resolution proof for $\Gamma$.
\end{proof}
We will prove the contrapositive of compactness theorem for the next 2 qn.
\section*{Question 16}
Assume that the Hilbert-style Proof System is sound and complete for propositional logic, Use these facts to establish that the Compactness Theorem holds true for propositional logic
\begin{proof}
    $\Gamma$ is unsat $\implies\Gamma \models \bot$, so by completeness $\Gamma \vdash \bot$ since proofs are finite, only finitely many assumptions from $\Gamma$ appear in the proof, yielding a finite $\Gamma_0 \subseteq_{\mathrm{fin}} \Gamma$ with $\Gamma_0 \vdash \bot$, 
    and by soundness $\Gamma_0 \models \bot$, hence $\Gamma_0$ is unsat
    \end{proof}
\section*{Question 17}
Assume that the Resolution Proof System is sound and complete for propositional logic. Use these facts to establish that the Compactness Theorem holds true for propositional logic
\begin{proof}
    If $\Gamma$ is an unsatidfiable set of clauses, by completeness there is a finite resolution proof for $\Gamma$, and since this proof uses only finitely many clauses from $\Gamma$, there is a finite subset $\Gamma_0 \subseteq_{\mathrm{fin}} \Gamma$ from which the same proof is a resolution proof. 
    By soundness, $\Gamma_0$ is unsat
    \end{proof}

\end{document}